\section{Finite ordered multigraphs}

\begin{definition}[Ordered directed multigraph]
An ordered directed multigraph is a tuple
\[
  G=(V,E,s,d,\prec)
\]
where $V=\{0,\ldots,n-1\}$ and $E=\{0,\ldots,m-1\}$ are finite,
$s,d:E\to V$ give the source and destination of each edge, and $\prec$
totally orders the outgoing edges of every vertex.  Parallel edges and
self-loops are permitted.
\end{definition}

The executable model represents $G$ as an adjacency-row sequence
$R_G=[R_0,\ldots,R_{n-1}]$.  Each $R_v$ is an ordered sequence of destination
identifiers.  The global edge identifier of the $i$th item in $R_v$ is
\[
  \mathsf{base}_G(v)+i,
  \qquad
  \mathsf{base}_G(v)=\sum_{u<v}|R_u|.
\]
This is precisely the logical information carried by a canonical CSR layout,
without making CSR part of the abstract interface.

\begin{definition}[Frontier]
A frontier is a finite sequence $F\in V^*$.  It is not a set.  Its order and
multiplicity are observable: if $v$ occurs twice, its outgoing row is traversed
twice.
\end{definition}

\begin{definition}[Edge context]
An edge context is a triple
\[
  c=(\mathsf{src}(c),\mathsf{dst}(c),\mathsf{eid}(c))\in V\times V\times E.
\]
For $v\in V$, let $X_G(v)$ be the sequence of contexts obtained from $R_v$ in
row order.  For $F=[v_0,\ldots,v_{k-1}]$, define
\[
  \trav_G(F)=X_G(v_0)\mathbin{+\!+}\cdots\mathbin{+\!+}X_G(v_{k-1}),
\]
where $+\!+$ denotes sequence concatenation.
\end{definition}

The regression-oracle model uses natural-number identifiers and assigns no
meaning to an invalid graph or to a frontier identifier outside $V$; its Lean
functions remain total for convenient execution.  The completeness model uses
the type \texttt{Fin $n$} for every vertex reference instead.  Valid sources,
destinations, and frontier entries are therefore enforced by construction in
the simulation theorems rather than carried as side conditions.
